\section{Discussion}
\label{sec:discussion}

\subsection{Contrast with the tetramer chain}

The main contrast with the tetramer chain follows from unit-cell parity.
Four spin-$1/2$ degrees of freedom can form an integer-spin or singlet cluster
sector, which is why the tetramer problem admits gapped tetramer and Haldane-like
phases separated by a gapless intermediate
region~\cite{li2026tetramer,gong2008tetramer}.
Five spins leave a half-integer degree of freedom unless translation symmetry is
broken or the ground state is degenerate, and the Lieb--Schultz--Mattis-type
constraints~\cite{lieb1961,ogata2019,affleck1987} then exclude a unique
symmetry-preserving gapped ground state under the stated assumptions.
At the tested positive-$\alpha$ points, the surviving doublet converts the
inter-cell exchange into an antiferromagnetic interaction between effective spins.
The finite-size observables accordingly point to a critical effective chain, and
the pentamer has to be classified on its own terms rather than through the
tetramer phase names and their triplon, quarton, or quinton labels.

The spectral comparison sharpens this distinction.
The pentamer's odd-parity excited doublet is dark at $q=\pi$, while higher
even-parity spin-$3/2$ sectors remain bright.
The finite-chain spectra contain low-energy weight and higher-energy structure,
but the local hierarchy alone does not identify the composition of those bands.
The multiplet analysis predicts allowed local channels, whereas a quantitative
assignment in the chain requires additional matrix-element or continuation
calculations.
A qualitatively similar separation between momentum-selective local weight and a
collective continuum has been discussed for other cluster and chain
magnets~\cite{keselman2020,gao2024srv3o9,bruognolo2016}, in line with the
interpretation given here.

The odd-cell constraint excludes a unique symmetric gapped ground state while
leaving open gapped symmetry-broken phases, degenerate ground states, and gapless
phases of other origin.
Among these, the finite-size evidence at the tested points selects the $c=1$
effective chain.
String-order and entanglement-spectrum
diagnostics~\cite{guwen2009,pollmann2010,pollmann2012,han2017bcfg} would be the
natural tools for probing hidden order in the same framework.

\subsection{Controlled interpretation near \texorpdfstring{$\alpha=0$}{alpha=0}}

Both sides of the coupling plane meet at $\alpha=0$, where the end bonds of each
cell vanish.
Sites 2--4 then form a trimer with a spin-$1/2$ ground doublet, while site 5 of
one cell and site 1 of the next are joined only by $J_3$ into a singlet dimer.
For $|\alpha|\ll\min(1,\beta)$ the weak end bonds let neighboring trimer doublets
interact through a virtual triplet excitation of the dimer, which costs an energy
of order $\beta$.
Once the end spins are bound into these dimers, the trimer doublets are the
low-energy degrees of freedom for either sign of $\alpha$, even though the isolated
cell, which lacks the inter-cell bond, has a quartet ground state for $\alpha<0$.
Projecting onto the trimer doublets and eliminating the dimer triplet gives
\begin{equation}
\Jeff=\frac{2\alpha^2}{9\beta}>0 .
\label{eq:near-zero}
\end{equation}
Exact diagonalization of an eight-spin trimer--dimer--trimer cluster at 47 points
supports this
local $\alpha^2/\beta$ scaling in its asymptotic domain, consistent with the
effective-spin logic established for trimer
chains~\cite{cheng2022trimer,cheng2024trimerfield}.
The finite-$\alpha$ controls at $(-0.10,0.25)$ and $(-0.175,0.25)$ deviate from the
leading term by $-36.1\%$ and $-53.9\%$, respectively.
These are signed deviations of the local approximation itself, defined as
$\epsilon=(\Delta_8/\alpha^2-2/(9\beta))/(2/(9\beta))$, where $\Delta_8$ is the
singlet--triplet gap of the eight-spin cluster, which for two coupled effective
spins equals $\Jeff$; the deviations grow as
$\alpha$ leaves the asymptotic domain, as expected for a leading-order result.
The decomposition thus supplies the mechanism near $\alpha=0$: it explains why the
effective exchange is antiferromagnetic and why it scales as $\alpha^2/\beta$.
It also explains why the near-zero strip is the most delicate region numerically:
the many-body energy scale vanishes quadratically as $\alpha\to0$, so finite-size
and solver effects compete with a very small exchange there.

\subsection{Meaning and boundary of the negative-\texorpdfstring{$\alpha$}{alpha} result}
\label{sec:negative-boundary}

The negative-$\alpha$ scan finds a change of the dominant finite-size momentum
across the sampled coupling plane at the retained solver settings.
In the thermodynamic limit such a reconstruction could sharpen into a Lifshitz
transition, spread into a broad incommensurate regime, or remain a crossover
driven by the competing-momentum mechanism known from frustrated and dimerized
chains~\cite{hida1992,chitra1995,hakobyan2001,white2008s1}.
The mismatch between the two diagnostics is compatible with a crossover but does
not distinguish these possibilities in the thermodynamic limit.
The local quartet analysis supplies a candidate origin for the deep-region value:
among the three inequivalent momenta folded from $Q=\pi$, the isolated
spin-$3/2$ cell has its largest shape-factor weight at $0.60\pi$ throughout the
sampled negative region, matching the measured deep-region $\qpeak$.
This is a comparison of folded replicas, not a claim that the isolated shape
factor has its global maximum there.
Whether the extended chain develops a singularity at that replica is a separate
many-body question, and the near-zero strip is especially sensitive to
initialization and fitting window.
The correspondence motivates a test: in the regime where the quartet projection
is controlled, larger-chain spectra should approach the relative replica weights
of Eq.~\eqref{eq:folded-response}.

The refinement calculations pass the stated acceptance gates, and the
$L=80$, $(\alpha,\beta)=(-0.10,0.25)$ endpoint is stable under
$k=8\rightarrow12$.
That local control does not exclude solver bias throughout the negative-region
scan; in particular, the $L=40$ and 60 members of the entropy size sequence still
use $k=3$.
Momentum reconstruction and finite-size flow are distinct observables, and at the
available sizes and coupling resolution they select different intervals.
The present evidence supports a finite-size momentum reconstruction.
A thermodynamic phase assignment requires a consistently controlled size sequence
and an appropriate bulk diagnostic; coincident signatures in the present
operational test would not alone establish such an assignment.

\subsection{Outlook}
\label{sec:remaining-tests}

The negative-$\alpha$ result sets up a sharp test of the phase question.
An observable with a controlled large-size singularity or crossing that tracks the
momentum reconstruction would test a thermodynamic boundary.
Figure~\ref{fig:fullq} now provides a negative-$\alpha$ dynamical example at
$(-0.5,0.5)$; a coupling scan and longer-time, larger-size spectra are needed to
connect that example to the static reconstruction intervals.
The five-momentum joint refinement should also be extended across the full grid,
with separate time-step and bond-dimension checks where quantitative line shapes
are required.

Two further extensions would connect the present results to the wider literature
on one-dimensional spin dynamics.
First, the momentum-resolved line shapes could be compared with the quasiparticle
interaction signatures and spinon hydrodynamics that have been analysed in
related chains~\cite{keselman2020,balents2020,wang2022hydro}, and with the
dynamical signatures used to distinguish symmetry-broken from liquid
phases~\cite{drescher2023,li2023dqcpsteer}.
Second, the entanglement analysis could be extended beyond the single-scalar
$c_{\mathrm{eff}}$ reported here to the full entanglement spectrum and its
dependence on microscopic anisotropy~\cite{kargarian2008}, and the momentum-space
continuum could be analysed with the same tools that have been applied to
honeycomb and square-lattice
antiferromagnets~\cite{gu2022honeycomb,wu2022j1j3}.

\subsection{Relation to experiment}

The local selection rule has a direct experimental consequence.
For an isolated reflection-symmetric cell, parity removes the lowest excited
doublet from the $q=\pi$ single-spin channel.
In a sufficiently weakly coupled pentamer material, this motivates a suppressed
lowest-doublet signal and stronger allowed spin-$3/2$ channels in inelastic
neutron scattering or single-spin $L$-edge RIXS.
Coupling-induced mixing and the relevant scattering operator must be checked
before extending the local selection rule quantitatively to a material.
Such a comparison follows the spectral hierarchy used to
identify composite excitations in trimer and tetramer
compounds~\cite{cheng2022trimer,li2025rixstrimer,li2026tetramer,prabhakar2025trimer}.
The prediction concerns the single-spin response, which is the channel measured
by neutrons and $L$-edge RIXS.
Because the Hamiltonian is SU(2) symmetric and the finite-chain ground state at
the calculated points is a spin singlet, $S^{zz}(q,\omega)$ fixes the unpolarized
neutron response up to a factor of three.
Two-particle channels probed by $K$-edge RIXS, and their relation to the
multi-spinon continuum observed in one-dimensional
cuprates~\cite{schlappa2018,ishii2025kedge,kumar2020,forte2011}, are a natural
extension of the same framework.
